\documentclass[11pt,a4paper]{article}
\usepackage[T1]{fontenc}
\usepackage[utf8]{inputenc}
\usepackage{amsmath,amssymb,xcolor}
\usepackage[margin=25mm]{geometry}
\usepackage[colorlinks=true,urlcolor=blue]{hyperref}
\definecolor{nexorablue}{HTML}{2459B5}
\setlength{\parindent}{0pt}
\setlength{\parskip}{7pt}
\setlength{\emergencystretch}{2em}
\renewcommand{\familydefault}{\sfdefault}
\begin{document}

{\small\bfseries\color{nexorablue} NEXORA CPGE}\par
{\small MP / KUT-SPE-07-02 / Corrigé}\par
{\LARGE\bfseries Obstruction à la diagonalisation}\par
\vspace{0.25em}\hrule\vspace{0.6em}

\section*{Multiplicité et noyau}
Le polynôme caractéristique est $(X-1)^2$. On a $(A-I)(x,y)^\top=(y,0)^\top$. Ainsi $E_1=\operatorname{Vect}((1,0)^\top)$.

\section*{Conclusion}
La seule valeur propre fournit une droite. Deux vecteurs propres ne peuvent donc constituer une base de dimension deux. La matrice est triangulaire mais non diagonalisable.

\bigskip
{\footnotesize Adaptation française par Nexora CPGE. Source : Ken Kuttler (Lyryx),
\href{https://math.libretexts.org/Bookshelves/Linear_Algebra/A_First_Course_in_Linear_Algebra_(Kuttler)/07\%253A_Spectral_Theory/7.02\%253A_Diagonalization}{\textit{A First Course in Linear Algebra}}, § 7.2, exemple 2.
\href{https://creativecommons.org/licenses/by/4.0/}{CC BY 4.0}.}
\end{document}
